% Insert in Section III-B, before the data-driven model comparison.
% Requires amsmath and graphicx. Merge ITU_References.bib into the manuscript bibliography.
% Copy the revised Fig04 PDF into the manuscript's Picture/ folder.
% Figure numbering is automatic; do not force the number with \renewcommand.

\subsubsection{Comparison with an ITU-based combined attenuation model}
\label{subsubsec:itu_combined_reference}

An ITU-based physical baseline combines gaseous, cloud, and rain
attenuation using the concurrent ERA5 inputs. Its hourly prediction is
\begin{equation}
\widehat{A}_{h}^{\mathrm{ITU}}
= A_{g,h}-A_{g,\mathrm{ref}}+A_{c,h}+A_{r,h},
\label{eq:itu_hourly_combination}
\end{equation}
where $A_{g,h}$, $A_{c,h}$, and $A_{r,h}$ denote the three path
attenuation components in decibels. The fixed gaseous reference
$A_{g,\mathrm{ref}}$ places the prediction on a relative-attenuation
scale without fitting an offset to the measured attenuation.

Gaseous attenuation follows the instantaneous oxygen and water-vapour
methods in Annex~2 of ITU-R P.676-13 \cite{ITU676_13}:
\begin{equation}
A_{g,h}
= \frac{\gamma_{o,h}H_{o,h}+K_{v,h}V_h}{\sin\theta},
\label{eq:itu_gaseous_hourly}
\end{equation}
where $\gamma_{o,h}$ is the oxygen specific attenuation in dB/km,
$H_{o,h}$ is its equivalent height in kilometres, $V_h$ is the
integrated water-vapour content in kg/m$^2$, and $K_{v,h}$ is the
water-vapour mass absorption coefficient in dB/(kg/m$^2$).
The angle $\theta$ is the link elevation.
The coefficients depend on frequency and the surface atmospheric
state. Water-vapour density is obtained from dewpoint and temperature
using ITU-R P.453-14 \cite{ITU453_14}.
For the local link, the calculation uses 19.13~GHz, circular
polarization, and the nominal link elevation of $42.969^{\circ}$.
The fixed reference uses the sea-level state of ITU-R P.835-6
\cite{ITU835_6}: 288.15~K, 1013.25~hPa, and 7.5~g/m$^3$ water-vapour
density. Integrating an exponential humidity profile with a 2-km
scale height gives a reference water-vapour column of 15~kg/m$^2$
and $A_{g,\mathrm{ref}}=0.2750$~dB.

Cloud attenuation is calculated from the ERA5 integrated cloud
liquid-water content $L_h$ using ITU-R P.840-9 \cite{ITU840_9}:
\begin{equation}
A_{c,h}=\frac{K_L(f)L_h}{\sin\theta},
\label{eq:itu_cloud_hourly}
\end{equation}
where $L_h$ is in kg/m$^2$ and $K_L(f)$ is the standard's
instantaneous cloud mass absorption coefficient in dB/(kg/m$^2$).
It includes the prescribed frequency correction to the liquid-water
coefficient evaluated at 273.75~K.

The rain component uses the specific-attenuation relation of
ITU-R P.838-3 \cite{ITU838_3} and a geometric rain path:
\begin{equation}
\begin{aligned}
\gamma_{r,h}&=k_R R_h^{\alpha_R},\\
\ell_{r,h}&=\frac{H_{0,h}+0.36}{\sin\theta},\\
A_{r,h}&=\gamma_{r,h}\ell_{r,h}.
\end{aligned}
\label{eq:itu_rain_hourly}
\end{equation}
Here, $\gamma_{r,h}$ is in dB/km and $\ell_{r,h}$ is in kilometres.
The hourly liquid-rain rate $R_h$ is in mm/h, obtained from
the time-aligned one-hour total precipitation minus snowfall water
equivalent; $k_R$ and $\alpha_R$ are the frequency- and
polarization-dependent coefficients. The nonnegative $H_{0,h}$ is
the ERA5 zero-degree height above ground in kilometres; unavailable
values use the fixed training-set median of 1.9724~km.
The 0.36-km increment follows ITU-R P.839-4 \cite{ITU839_4}.
Using an hourly zero-degree height and rain rate makes this an
hourly geometric-path approximation, rather than an annual
exceedance-probability calculation. The combined prediction does not
include a separate scintillation component.

\begin{figure*}[t]
\centering
\includegraphics[width=\textwidth]{Picture/Fig04_All_Hourly_Observations_And_ITU_Test.pdf}
\caption{Hourly attenuation observations and the ITU-based test
comparison. (a) All 3434 local measured hours across the training,
validation, and test partitions, with background colours denoting
weather categories. (b) The measured attenuation and ITU-based
combined predictions for the 684 test hours with valid physical
inputs. Each panel has its own chronological sample index, with
calendar gaps omitted. Panel (b) retains the enlarged vertical
range; triangles mark 15 observations above its upper limit.
Panel (a) retains the complete measured attenuation range.
The reported scores use all 684 pairs without clipping.}
\label{fig:itu_complete_hourly_sequence}
\end{figure*}

Figure~\ref{fig:itu_complete_hourly_sequence}(a) shows all 3434
local hourly attenuation observations with their weather labels.
Figure~\ref{fig:itu_complete_hourly_sequence}(b) compares the
measured attenuation and the ITU-based combined predictions on
the common 684-hour test subset. Three of the original 687 test hours have a negative
total-precipitation-minus-snowfall difference and are excluded
from this physical-baseline comparison. The combined model
obtains an RMSE of 1.0090~dB, an MAE of 0.4317~dB, and
$R^2=0.0281$. Its mean prediction-minus-observation bias is
$-0.1193$~dB. The predicted sequence has a standard deviation of
0.3304~dB, compared with 1.0234~dB for the observations, indicating
that this hourly implementation reproduces substantially less
variability than measured on the link.

The discrepancy is not a uniform offset. For the 51 rainy test
hours, the RMSE is 1.3559~dB and the mean bias is $+0.2951$~dB,
although the larger observed attenuation excursions remain
underrepresented. The RMSE is 0.9747~dB for the 156 dry-overcast
hours. These results establish a physical reference for the
following data-driven comparisons: the relevant question is how
well concurrent atmospheric information can reproduce the
measured link-specific hourly variations beyond this prescribed
physical approximation.
